% =====================================================================
% additional.tex  --  BODY ONLY. Do not add \documentclass or \begin{document}.
% Use from supp.tex (inside the document) with:   % =====================================================================
% additional.tex  --  BODY ONLY. Do not add \documentclass or \begin{document}.
% Use from supp.tex (inside the document) with:   % =====================================================================
% additional.tex  --  BODY ONLY. Do not add \documentclass or \begin{document}.
% Use from supp.tex (inside the document) with:   % =====================================================================
% additional.tex  --  BODY ONLY. Do not add \documentclass or \begin{document}.
% Use from supp.tex (inside the document) with:   \input{additional}
% Needs amsmath + amssymb (already loaded by CVPR preambles).
% Add the two entries in refs_additions.bib to refs.bib.
% =====================================================================
\makeatletter
\@ifundefined{c@suppprop}{\newcounter{suppprop}}{}
\renewcommand{\thesuppprop}{S\arabic{suppprop}}
\@ifundefined{proposition}{%
  \newenvironment{proposition}[1][]{\refstepcounter{suppprop}\par\medskip\noindent
    \textbf{Proposition~\thesuppprop}\if\relax\detokenize{#1}\relax\else\ \textbf{(#1)}\fi\textbf{.}\ \itshape}%
   {\par\medskip}}{}
\makeatother
\providecommand{\auroc}{\textsc{auroc}}
\providecommand{\cur}{\textsc{cur}}
\providecommand{\E}{\mathbb{E}}

\section{Why the method works: logical and mathematical justification}
\label{sec:supp-why}

This section states, under explicit assumptions, why native confidence cannot
detect missing support, why cross-view tokens can, why a head trained on its
backbone's own reprojection labels can match a ground-truth head and beat its
teacher, and why no detector of the inputs can catch an invention every view
agrees on. It ends with a check that the reported \auroc{} and \cur{} values
agree with each other.

\paragraph{Notation.}
View $i$ of image set $\mathcal{I}=\{I_1,\dots,I_N\}$; pixel $p$; support label
$Y\!\in\!\{0,1\}$ ($1$ for supported, S; $0$ for unsupported, U); depth error
$e=|\hat D_i(p)-D_i(p)|/D_i(p)$. Per-image tokens
$z^{0}_i=E(I_i)$ depend on $I_i$ only; cross-view tokens
$z_i=A(\mathcal{I})$ depend on all views. A score $s$ is any real function;
$\auroc(s)=P(s_S>s_U)+\tfrac12P(s_S=s_U)$, and $\cur_\tau(s)=P(s_U>\tau)$ with
$\tau$ chosen so that $P(s_S>\tau)=0.9$.

% ---------------------------------------------------------------------------
\subsection{Native confidence predicts error, not support}

DUSt3R and its descendants train confidence with a confidence-weighted
regression loss~\cite{wang2024dust3r}
\begin{equation}
  \mathcal{L}(C)=C\,e-\alpha\log C,\qquad \alpha>0,
  \label{eq:confloss}
\end{equation}
applied only where ground-truth depth exists. For backbones whose exact loss
differs, the argument below holds for any loss whose minimiser is a
decreasing function of expected error.

\begin{proposition}[Loss-optimal confidence]\label{prop:conf}
For a head with input $z$, the minimiser of $\E[\mathcal{L}(C)\mid z]$ is
\begin{equation}
  C^\star(z)=\frac{\alpha}{m(z)},\qquad m(z)=\E[e\mid z].
  \label{eq:copt}
\end{equation}
Consequently $C^\star$ carries information about $Y$ only through $m(z)$.
If $e\perp Y\mid z$, then $C^\star$ and $Y$ are conditionally independent
given $m(z)$, and $\auroc(C^\star)$ for S versus U is the \auroc{} of expected
error for separating S from U.
\end{proposition}
\par\noindent\emph{Proof.}~
$\partial_C\E[\mathcal{L}\mid z]=m(z)-\alpha/C=0$ gives Eq.~\eqref{eq:copt};
the second derivative $\alpha/C^2>0$ makes this the minimum. $C^\star$ is a
function of $m(z)$ alone.
\hfill$\square$\par

\noindent\emph{Reading.} The objective contains no term that depends on $Y$, so
nothing pushes the head to encode support even when $z$ contains it. Invented
depth is accurate (\textsc{AbsRel}$_U$ 0.011 to 0.022, main paper, Sec.~4), so
$m(z)$ is nearly the same on S and U, and Eq.~\eqref{eq:copt} predicts high
confidence on invented geometry. This is the leak in main-paper Table~1.
Proposition~\ref{prop:conf} also explains main-paper Gate~11: filtering for
accuracy is an error-ranking task, which $C^\star$ already performs, so adding
a support score does not raise fused-cloud F1. Support and accuracy are
different targets and need different scores.

% ---------------------------------------------------------------------------
\subsection{The matched split forces chance for any single-view score}

\begin{proposition}[Exchangeability]\label{prop:exch}
Suppose that, in the scored view $i$, the pixels of patch A (U) and patch B
(S$_m$) are exchangeable: the joint law of $(I_i,p)$ for $p\!\in\!A$ equals
that for $p\!\in\!B$ (matched area, depth, texture and fill, with roles
assigned independently of view-$i$ appearance). Then every score
$s=f(I_i,p)$, including any function of $z^0_i$, satisfies
\begin{equation}
  \auroc_m(s)=\tfrac12,\qquad \cur_m(s)=0.90 .
  \label{eq:chance}
\end{equation}
\end{proposition}
\par\noindent\emph{Proof.}~
Exchangeability gives $s_A\overset{d}{=}s_B$, so for independent draws
$P(s_B>s_A)=P(s_A>s_B)$ and $\auroc_m=\tfrac12$. With identical laws, the
threshold that passes 90\% of B passes 90\% of A.
\hfill$\square$\par

\begin{proposition}[Information ordering]\label{prop:order}
Let $\mathcal{F}_0=\{f(I_i,p)\}$ and $\mathcal{F}=\{f(\mathcal{I},p)\}$. Since
$\mathcal{F}_0\subset\mathcal{F}$,
$\sup_{\mathcal{F}}\auroc\ge\sup_{\mathcal{F}_0}\auroc$, and the supremum over
each class is attained by the posterior $\eta(x)=P(Y\!=\!1\mid x)$ of its
input~\cite{menon2015corrupted}. Under the matched split, any
$\auroc_m>\tfrac12$ therefore measures information that comes from views
other than $i$.
\end{proposition}

\noindent\emph{Reading.} Eq.~\eqref{eq:chance} gives the reference values for
the shaded columns of main-paper Table~1. Native confidence scores
$\auroc_m$ 0.43 to 0.55 and $\cur_m$ 0.80 to 0.87, close to chance: it acts as
a single-view score even though the backbone's tokens have cross-view inputs,
as Proposition~\ref{prop:conf} predicts. The per-image control is a member of
$\mathcal{F}_0$ by construction, so Eq.~\eqref{eq:chance} explains why it stays
at 0.56 to 0.61. The residual above $\tfrac12$ is the cost of approximate
matching (image-border and placement cues), which the pre-registered 0.60 bound
caps. Proposition~\ref{prop:order} explains why the control is decisive: the
support head and the control share architecture, labels and training, and
differ only in whether their input is in $\mathcal{F}$ or $\mathcal{F}_0$.

% ---------------------------------------------------------------------------
\subsection{Why cross-view tokens carry support}

Support is defined over the view set: $Y(p)=1$ iff the ground-truth point at
$p$ is visible in some view $j\neq i$. It is a property of $\mathcal{I}$, not
of $I_i$. Cross-view attention computes, for each token, a weighting over
tokens in the other views; when the surface at $p$ appears in view $j$, the
token has a matching partner there, and when it does not, it has none. The
backbone must compute this correspondence to triangulate, so the information
is available whether or not any head reads it. A linear probe reaches
\auroc{} 0.99 (main paper, Figure~8). This shows the information is linearly
decodable, which is why a head of 1.7 to 3.3\,M parameters is enough and why a
probe on all four tapped layers is within 0.015 of any head we trained.

\paragraph{Evidence grows with views.}
Model each other view $j$ as contributing an independent log-likelihood-ratio
term $\ell_j$ for $Y$, with mean $\pm\mu/2$ and variance $\sigma^2$. Summing
$k$ views gives separation $d'_k=\sqrt{k}\,\mu/\sigma$, and under a binormal
score model
\begin{equation}
  \auroc_k=\Phi\!\big(\sqrt{k}\,d'_1/\sqrt2\big).
  \label{eq:views}
\end{equation}
From a one-view value of 0.68 to 0.75, Eq.~\eqref{eq:views} predicts 0.75 to 0.83
for two views and 0.79 to 0.88 for three. The measured values are 0.80 to 0.92
for three or more views (main paper, Sec.~5). The per-view term $\mu$ scales with
parallax, which triangulation needs, so a single low-parallax view gives
little evidence. That is consistent with main-paper Gate~9: the one-view fine-tune
stops at 0.71, short of 0.75, because there is little evidence to read, not
because the head is too small.

% ---------------------------------------------------------------------------
\subsection{Learning from the backbone's own reprojection labels}

The label-free head trains on $\tilde Y$, the reprojection test on the
backbone's own depth and cameras. Treat $\tilde Y$ as a corrupted copy of $Y$
with class-conditional flip rates $\rho_{+}=P(\tilde Y\!=\!0\mid Y\!=\!1)$ and
$\rho_{-}=P(\tilde Y\!=\!1\mid Y\!=\!0)$.

\begin{proposition}[Rank preservation~\cite{menon2015corrupted}]\label{prop:noise}
With $\eta(z)=P(Y\!=\!1\mid z)$ and $\tilde\eta(z)=P(\tilde Y\!=\!1\mid z)$,
\begin{equation}
  \tilde\eta(z)=\rho_{-}+(1-\rho_{+}-\rho_{-})\,\eta(z).
  \label{eq:noisy}
\end{equation}
If $\rho_{+}+\rho_{-}<1$, then $\tilde\eta$ is a strictly increasing affine
function of $\eta$ and $\auroc(\tilde\eta)=\auroc(\eta)$. A head that estimates
$\tilde\eta$ by binary cross-entropy has, in the limit, the same \auroc{} as a
head trained on ground truth.
\end{proposition}
\par\noindent\emph{Proof.}~
$\tilde\eta=P(\tilde Y\!=\!1\mid Y\!=\!1)\eta+P(\tilde Y\!=\!1\mid Y\!=\!0)(1-\eta)
=(1-\rho_{+})\eta+\rho_{-}(1-\eta)$. \auroc{} depends only on the ordering of
scores, which a strictly increasing map preserves. Cross-entropy is a proper
loss, so its population minimiser is $\tilde\eta$.
\hfill$\square$\par

\noindent\emph{Why the head can beat its teacher.} The teacher scores each
pixel with its own reprojection, $t=g(\eta(z))+\xi$, where $\xi$ is
per-pixel depth and pose noise that the tokens do not predict. The head
estimates $\E[\tilde Y\mid z]$ and so averages $\xi$ out. Under a binormal
model with signal variance $\sigma^2$ and noise variance $\sigma_\xi^2$, the
teacher's separation drops to $d'/\sqrt{1+\sigma_\xi^2/\sigma^2}$, while the
head keeps $d'$. This matches main-paper Gates~3 and~4: on DTU the label-free
head reaches 0.99 against 0.94 for reprojection.

\noindent\emph{When it fails.} Eq.~\eqref{eq:noisy} needs
$\rho_{+}+\rho_{-}<1$, and the gap shrinks as $(1-\rho_{+}-\rho_{-})$ falls.
In practice $\rho$ also depends on $z$, which can reorder pixels. Both effects
are largest where the teacher's geometry is poor. This is why adapting
VGGT label-free on CO3Dv2 loses 0.09, and why the DTU-trained head is the safe
default (main paper, Sec.~6).

% ---------------------------------------------------------------------------
\subsection{Why labels from fills need the matched control}

Let $F$ be any appearance cue, such as the fill. If training makes $F$ predict
$\tilde Y$, the loss-optimal head uses $F$; this is what happens with the
benchmark's fills (\auroc{} 1.000, with the per-image control learning the
same shortcut). Matched labels make $F$ uninformative by design:
$P(F\mid\tilde Y\!=\!1)=P(F\mid\tilde Y\!=\!0)$, so $I(\tilde Y;F)=0$. The
population minimiser then gains nothing from $F$ and must use what still
predicts $\tilde Y$, which by Proposition~\ref{prop:order} is cross-view
evidence. The per-image control is the test that this worked: it must stay at
chance under the same labels.

% ---------------------------------------------------------------------------
\subsection{The limit: consistency is not observation}

\begin{proposition}[Indistinguishability]\label{prop:limit}
Let world $W_1$ contain a surface $X$ that no view sees, which every view shows
as a patch $P$, and let $W_2$ contain a real surface that looks like $P$ in
every view. If both worlds produce the same images $\mathcal{I}$, then every
score $s=f(\mathcal{I},p)$, and every test built from the backbone's outputs
(which are themselves functions of $\mathcal{I}$), gives the same value in both.
No such score can label the patch pixels U in $W_1$ and S in $W_2$.
\end{proposition}
\par\noindent\emph{Proof.}~
$s$ is a function of $\mathcal{I}$ and $p$, which are equal in the two worlds.
\hfill$\square$\par

\noindent\emph{Reading.} This is main-paper Figure~13 and the first limitation
in Sec.~7. Reprojection, model disagreement and the support head are all
functions of $\mathcal{I}$, so none of them separates a patch that every view
shows. Refilling the hole does separate it, because refilling changes
$\mathcal{I}$: it is an intervention, not a function of the inputs.
Consequently the head certifies multi-view agreement, not truth. Reprojection
consistency is necessary for support, not sufficient, which is why
reprojection is at chance on the patches (0.49, 0.48).

% ---------------------------------------------------------------------------
\subsection{Consistency check between \auroc{} and \cur}

If scores on S and U are Gaussian with equal variance and separation $d'$,
then $\auroc=\Phi(d'/\sqrt2)$, and with $z_{0.9}=\Phi^{-1}(0.9)=1.2816$,
\begin{equation}
  \cur=\Phi\!\big(z_{0.9}-\sqrt2\,\Phi^{-1}(\auroc)\big).
  \label{eq:binormal}
\end{equation}
Eq.~\eqref{eq:binormal} predicts each table's \cur{} from its \auroc{} alone
(Table~\ref{tab:binormal}). For the support head it fits within 0.02 on all
four datasets, and for native confidence on DTU within 0.01. The two metrics
are therefore consistent with each other, not two separate results. The fit is
worst for reprojection on Photo Tourism (0.16), where the score is a count of
views and the Gaussian model does not hold. On the matched split the model
predicts 0.87 to 0.94, above the measured 0.80 to 0.87, because the two
patches' score distributions differ in spread even when their medians agree.

\begin{table}[!t]
\centering\small
\setlength{\tabcolsep}{4pt}
\begin{tabular}{@{}llccc@{}}
\toprule
detector & data & \auroc & \cur{} meas. & \cur{} pred. \\
\midrule
native conf.   & DTU     & 0.82 & 0.49 & 0.50 \\
native conf.   & CO3Dv2  & 0.67 & 0.80 & 0.75 \\
native conf.   & ETH3D   & 0.76 & 0.50 & 0.61 \\
native conf.   & Photot. & 0.82 & 0.46 & 0.50 \\
reprojection   & DTU     & 0.94 & 0.13 & 0.18 \\
reprojection   & ETH3D   & 0.97 & 0.07 & 0.08 \\
reprojection   & Photot. & 0.82 & 0.34 & 0.50 \\
head, label-free & DTU   & 0.99 & 0.03 & 0.02 \\
head, label-free & CO3Dv2& 0.92 & 0.23 & 0.24 \\
head, label-free & ETH3D & 0.97 & 0.09 & 0.08 \\
head, label-free & Photot.& 0.93 & 0.19 & 0.21 \\
\midrule
chance (Prop.~\ref{prop:exch}) & any & 0.50 & --- & 0.90 \\
\bottomrule
\end{tabular}
\caption{Measured \cur{} against the binormal prediction of
Eq.~\eqref{eq:binormal} from \auroc{} alone (values from main-paper
Table~2, averaged over backbones).}
\label{tab:binormal}
\end{table}
% Needs amsmath + amssymb (already loaded by CVPR preambles).
% Add the two entries in refs_additions.bib to refs.bib.
% =====================================================================
\makeatletter
\@ifundefined{c@suppprop}{\newcounter{suppprop}}{}
\renewcommand{\thesuppprop}{S\arabic{suppprop}}
\@ifundefined{proposition}{%
  \newenvironment{proposition}[1][]{\refstepcounter{suppprop}\par\medskip\noindent
    \textbf{Proposition~\thesuppprop}\if\relax\detokenize{#1}\relax\else\ \textbf{(#1)}\fi\textbf{.}\ \itshape}%
   {\par\medskip}}{}
\makeatother
\providecommand{\auroc}{\textsc{auroc}}
\providecommand{\cur}{\textsc{cur}}
\providecommand{\E}{\mathbb{E}}

\section{Why the method works: logical and mathematical justification}
\label{sec:supp-why}

This section states, under explicit assumptions, why native confidence cannot
detect missing support, why cross-view tokens can, why a head trained on its
backbone's own reprojection labels can match a ground-truth head and beat its
teacher, and why no detector of the inputs can catch an invention every view
agrees on. It ends with a check that the reported \auroc{} and \cur{} values
agree with each other.

\paragraph{Notation.}
View $i$ of image set $\mathcal{I}=\{I_1,\dots,I_N\}$; pixel $p$; support label
$Y\!\in\!\{0,1\}$ ($1$ for supported, S; $0$ for unsupported, U); depth error
$e=|\hat D_i(p)-D_i(p)|/D_i(p)$. Per-image tokens
$z^{0}_i=E(I_i)$ depend on $I_i$ only; cross-view tokens
$z_i=A(\mathcal{I})$ depend on all views. A score $s$ is any real function;
$\auroc(s)=P(s_S>s_U)+\tfrac12P(s_S=s_U)$, and $\cur_\tau(s)=P(s_U>\tau)$ with
$\tau$ chosen so that $P(s_S>\tau)=0.9$.

% ---------------------------------------------------------------------------
\subsection{Native confidence predicts error, not support}

DUSt3R and its descendants train confidence with a confidence-weighted
regression loss~\cite{wang2024dust3r}
\begin{equation}
  \mathcal{L}(C)=C\,e-\alpha\log C,\qquad \alpha>0,
  \label{eq:confloss}
\end{equation}
applied only where ground-truth depth exists. For backbones whose exact loss
differs, the argument below holds for any loss whose minimiser is a
decreasing function of expected error.

\begin{proposition}[Loss-optimal confidence]\label{prop:conf}
For a head with input $z$, the minimiser of $\E[\mathcal{L}(C)\mid z]$ is
\begin{equation}
  C^\star(z)=\frac{\alpha}{m(z)},\qquad m(z)=\E[e\mid z].
  \label{eq:copt}
\end{equation}
Consequently $C^\star$ carries information about $Y$ only through $m(z)$.
If $e\perp Y\mid z$, then $C^\star$ and $Y$ are conditionally independent
given $m(z)$, and $\auroc(C^\star)$ for S versus U is the \auroc{} of expected
error for separating S from U.
\end{proposition}
\par\noindent\emph{Proof.}~
$\partial_C\E[\mathcal{L}\mid z]=m(z)-\alpha/C=0$ gives Eq.~\eqref{eq:copt};
the second derivative $\alpha/C^2>0$ makes this the minimum. $C^\star$ is a
function of $m(z)$ alone.
\hfill$\square$\par

\noindent\emph{Reading.} The objective contains no term that depends on $Y$, so
nothing pushes the head to encode support even when $z$ contains it. Invented
depth is accurate (\textsc{AbsRel}$_U$ 0.011 to 0.022, main paper, Sec.~4), so
$m(z)$ is nearly the same on S and U, and Eq.~\eqref{eq:copt} predicts high
confidence on invented geometry. This is the leak in main-paper Table~1.
Proposition~\ref{prop:conf} also explains main-paper Gate~11: filtering for
accuracy is an error-ranking task, which $C^\star$ already performs, so adding
a support score does not raise fused-cloud F1. Support and accuracy are
different targets and need different scores.

% ---------------------------------------------------------------------------
\subsection{The matched split forces chance for any single-view score}

\begin{proposition}[Exchangeability]\label{prop:exch}
Suppose that, in the scored view $i$, the pixels of patch A (U) and patch B
(S$_m$) are exchangeable: the joint law of $(I_i,p)$ for $p\!\in\!A$ equals
that for $p\!\in\!B$ (matched area, depth, texture and fill, with roles
assigned independently of view-$i$ appearance). Then every score
$s=f(I_i,p)$, including any function of $z^0_i$, satisfies
\begin{equation}
  \auroc_m(s)=\tfrac12,\qquad \cur_m(s)=0.90 .
  \label{eq:chance}
\end{equation}
\end{proposition}
\par\noindent\emph{Proof.}~
Exchangeability gives $s_A\overset{d}{=}s_B$, so for independent draws
$P(s_B>s_A)=P(s_A>s_B)$ and $\auroc_m=\tfrac12$. With identical laws, the
threshold that passes 90\% of B passes 90\% of A.
\hfill$\square$\par

\begin{proposition}[Information ordering]\label{prop:order}
Let $\mathcal{F}_0=\{f(I_i,p)\}$ and $\mathcal{F}=\{f(\mathcal{I},p)\}$. Since
$\mathcal{F}_0\subset\mathcal{F}$,
$\sup_{\mathcal{F}}\auroc\ge\sup_{\mathcal{F}_0}\auroc$, and the supremum over
each class is attained by the posterior $\eta(x)=P(Y\!=\!1\mid x)$ of its
input~\cite{menon2015corrupted}. Under the matched split, any
$\auroc_m>\tfrac12$ therefore measures information that comes from views
other than $i$.
\end{proposition}

\noindent\emph{Reading.} Eq.~\eqref{eq:chance} gives the reference values for
the shaded columns of main-paper Table~1. Native confidence scores
$\auroc_m$ 0.43 to 0.55 and $\cur_m$ 0.80 to 0.87, close to chance: it acts as
a single-view score even though the backbone's tokens have cross-view inputs,
as Proposition~\ref{prop:conf} predicts. The per-image control is a member of
$\mathcal{F}_0$ by construction, so Eq.~\eqref{eq:chance} explains why it stays
at 0.56 to 0.61. The residual above $\tfrac12$ is the cost of approximate
matching (image-border and placement cues), which the pre-registered 0.60 bound
caps. Proposition~\ref{prop:order} explains why the control is decisive: the
support head and the control share architecture, labels and training, and
differ only in whether their input is in $\mathcal{F}$ or $\mathcal{F}_0$.

% ---------------------------------------------------------------------------
\subsection{Why cross-view tokens carry support}

Support is defined over the view set: $Y(p)=1$ iff the ground-truth point at
$p$ is visible in some view $j\neq i$. It is a property of $\mathcal{I}$, not
of $I_i$. Cross-view attention computes, for each token, a weighting over
tokens in the other views; when the surface at $p$ appears in view $j$, the
token has a matching partner there, and when it does not, it has none. The
backbone must compute this correspondence to triangulate, so the information
is available whether or not any head reads it. A linear probe reaches
\auroc{} 0.99 (main paper, Figure~8). This shows the information is linearly
decodable, which is why a head of 1.7 to 3.3\,M parameters is enough and why a
probe on all four tapped layers is within 0.015 of any head we trained.

\paragraph{Evidence grows with views.}
Model each other view $j$ as contributing an independent log-likelihood-ratio
term $\ell_j$ for $Y$, with mean $\pm\mu/2$ and variance $\sigma^2$. Summing
$k$ views gives separation $d'_k=\sqrt{k}\,\mu/\sigma$, and under a binormal
score model
\begin{equation}
  \auroc_k=\Phi\!\big(\sqrt{k}\,d'_1/\sqrt2\big).
  \label{eq:views}
\end{equation}
From a one-view value of 0.68 to 0.75, Eq.~\eqref{eq:views} predicts 0.75 to 0.83
for two views and 0.79 to 0.88 for three. The measured values are 0.80 to 0.92
for three or more views (main paper, Sec.~5). The per-view term $\mu$ scales with
parallax, which triangulation needs, so a single low-parallax view gives
little evidence. That is consistent with main-paper Gate~9: the one-view fine-tune
stops at 0.71, short of 0.75, because there is little evidence to read, not
because the head is too small.

% ---------------------------------------------------------------------------
\subsection{Learning from the backbone's own reprojection labels}

The label-free head trains on $\tilde Y$, the reprojection test on the
backbone's own depth and cameras. Treat $\tilde Y$ as a corrupted copy of $Y$
with class-conditional flip rates $\rho_{+}=P(\tilde Y\!=\!0\mid Y\!=\!1)$ and
$\rho_{-}=P(\tilde Y\!=\!1\mid Y\!=\!0)$.

\begin{proposition}[Rank preservation~\cite{menon2015corrupted}]\label{prop:noise}
With $\eta(z)=P(Y\!=\!1\mid z)$ and $\tilde\eta(z)=P(\tilde Y\!=\!1\mid z)$,
\begin{equation}
  \tilde\eta(z)=\rho_{-}+(1-\rho_{+}-\rho_{-})\,\eta(z).
  \label{eq:noisy}
\end{equation}
If $\rho_{+}+\rho_{-}<1$, then $\tilde\eta$ is a strictly increasing affine
function of $\eta$ and $\auroc(\tilde\eta)=\auroc(\eta)$. A head that estimates
$\tilde\eta$ by binary cross-entropy has, in the limit, the same \auroc{} as a
head trained on ground truth.
\end{proposition}
\par\noindent\emph{Proof.}~
$\tilde\eta=P(\tilde Y\!=\!1\mid Y\!=\!1)\eta+P(\tilde Y\!=\!1\mid Y\!=\!0)(1-\eta)
=(1-\rho_{+})\eta+\rho_{-}(1-\eta)$. \auroc{} depends only on the ordering of
scores, which a strictly increasing map preserves. Cross-entropy is a proper
loss, so its population minimiser is $\tilde\eta$.
\hfill$\square$\par

\noindent\emph{Why the head can beat its teacher.} The teacher scores each
pixel with its own reprojection, $t=g(\eta(z))+\xi$, where $\xi$ is
per-pixel depth and pose noise that the tokens do not predict. The head
estimates $\E[\tilde Y\mid z]$ and so averages $\xi$ out. Under a binormal
model with signal variance $\sigma^2$ and noise variance $\sigma_\xi^2$, the
teacher's separation drops to $d'/\sqrt{1+\sigma_\xi^2/\sigma^2}$, while the
head keeps $d'$. This matches main-paper Gates~3 and~4: on DTU the label-free
head reaches 0.99 against 0.94 for reprojection.

\noindent\emph{When it fails.} Eq.~\eqref{eq:noisy} needs
$\rho_{+}+\rho_{-}<1$, and the gap shrinks as $(1-\rho_{+}-\rho_{-})$ falls.
In practice $\rho$ also depends on $z$, which can reorder pixels. Both effects
are largest where the teacher's geometry is poor. This is why adapting
VGGT label-free on CO3Dv2 loses 0.09, and why the DTU-trained head is the safe
default (main paper, Sec.~6).

% ---------------------------------------------------------------------------
\subsection{Why labels from fills need the matched control}

Let $F$ be any appearance cue, such as the fill. If training makes $F$ predict
$\tilde Y$, the loss-optimal head uses $F$; this is what happens with the
benchmark's fills (\auroc{} 1.000, with the per-image control learning the
same shortcut). Matched labels make $F$ uninformative by design:
$P(F\mid\tilde Y\!=\!1)=P(F\mid\tilde Y\!=\!0)$, so $I(\tilde Y;F)=0$. The
population minimiser then gains nothing from $F$ and must use what still
predicts $\tilde Y$, which by Proposition~\ref{prop:order} is cross-view
evidence. The per-image control is the test that this worked: it must stay at
chance under the same labels.

% ---------------------------------------------------------------------------
\subsection{The limit: consistency is not observation}

\begin{proposition}[Indistinguishability]\label{prop:limit}
Let world $W_1$ contain a surface $X$ that no view sees, which every view shows
as a patch $P$, and let $W_2$ contain a real surface that looks like $P$ in
every view. If both worlds produce the same images $\mathcal{I}$, then every
score $s=f(\mathcal{I},p)$, and every test built from the backbone's outputs
(which are themselves functions of $\mathcal{I}$), gives the same value in both.
No such score can label the patch pixels U in $W_1$ and S in $W_2$.
\end{proposition}
\par\noindent\emph{Proof.}~
$s$ is a function of $\mathcal{I}$ and $p$, which are equal in the two worlds.
\hfill$\square$\par

\noindent\emph{Reading.} This is main-paper Figure~13 and the first limitation
in Sec.~7. Reprojection, model disagreement and the support head are all
functions of $\mathcal{I}$, so none of them separates a patch that every view
shows. Refilling the hole does separate it, because refilling changes
$\mathcal{I}$: it is an intervention, not a function of the inputs.
Consequently the head certifies multi-view agreement, not truth. Reprojection
consistency is necessary for support, not sufficient, which is why
reprojection is at chance on the patches (0.49, 0.48).

% ---------------------------------------------------------------------------
\subsection{Consistency check between \auroc{} and \cur}

If scores on S and U are Gaussian with equal variance and separation $d'$,
then $\auroc=\Phi(d'/\sqrt2)$, and with $z_{0.9}=\Phi^{-1}(0.9)=1.2816$,
\begin{equation}
  \cur=\Phi\!\big(z_{0.9}-\sqrt2\,\Phi^{-1}(\auroc)\big).
  \label{eq:binormal}
\end{equation}
Eq.~\eqref{eq:binormal} predicts each table's \cur{} from its \auroc{} alone
(Table~\ref{tab:binormal}). For the support head it fits within 0.02 on all
four datasets, and for native confidence on DTU within 0.01. The two metrics
are therefore consistent with each other, not two separate results. The fit is
worst for reprojection on Photo Tourism (0.16), where the score is a count of
views and the Gaussian model does not hold. On the matched split the model
predicts 0.87 to 0.94, above the measured 0.80 to 0.87, because the two
patches' score distributions differ in spread even when their medians agree.

\begin{table}[!t]
\centering\small
\setlength{\tabcolsep}{4pt}
\begin{tabular}{@{}llccc@{}}
\toprule
detector & data & \auroc & \cur{} meas. & \cur{} pred. \\
\midrule
native conf.   & DTU     & 0.82 & 0.49 & 0.50 \\
native conf.   & CO3Dv2  & 0.67 & 0.80 & 0.75 \\
native conf.   & ETH3D   & 0.76 & 0.50 & 0.61 \\
native conf.   & Photot. & 0.82 & 0.46 & 0.50 \\
reprojection   & DTU     & 0.94 & 0.13 & 0.18 \\
reprojection   & ETH3D   & 0.97 & 0.07 & 0.08 \\
reprojection   & Photot. & 0.82 & 0.34 & 0.50 \\
head, label-free & DTU   & 0.99 & 0.03 & 0.02 \\
head, label-free & CO3Dv2& 0.92 & 0.23 & 0.24 \\
head, label-free & ETH3D & 0.97 & 0.09 & 0.08 \\
head, label-free & Photot.& 0.93 & 0.19 & 0.21 \\
\midrule
chance (Prop.~\ref{prop:exch}) & any & 0.50 & --- & 0.90 \\
\bottomrule
\end{tabular}
\caption{Measured \cur{} against the binormal prediction of
Eq.~\eqref{eq:binormal} from \auroc{} alone (values from main-paper
Table~2, averaged over backbones).}
\label{tab:binormal}
\end{table}
% Needs amsmath + amssymb (already loaded by CVPR preambles).
% Add the two entries in refs_additions.bib to refs.bib.
% =====================================================================
\makeatletter
\@ifundefined{c@suppprop}{\newcounter{suppprop}}{}
\renewcommand{\thesuppprop}{S\arabic{suppprop}}
\@ifundefined{proposition}{%
  \newenvironment{proposition}[1][]{\refstepcounter{suppprop}\par\medskip\noindent
    \textbf{Proposition~\thesuppprop}\if\relax\detokenize{#1}\relax\else\ \textbf{(#1)}\fi\textbf{.}\ \itshape}%
   {\par\medskip}}{}
\makeatother
\providecommand{\auroc}{\textsc{auroc}}
\providecommand{\cur}{\textsc{cur}}
\providecommand{\E}{\mathbb{E}}

\section{Why the method works: logical and mathematical justification}
\label{sec:supp-why}

This section states, under explicit assumptions, why native confidence cannot
detect missing support, why cross-view tokens can, why a head trained on its
backbone's own reprojection labels can match a ground-truth head and beat its
teacher, and why no detector of the inputs can catch an invention every view
agrees on. It ends with a check that the reported \auroc{} and \cur{} values
agree with each other.

\paragraph{Notation.}
View $i$ of image set $\mathcal{I}=\{I_1,\dots,I_N\}$; pixel $p$; support label
$Y\!\in\!\{0,1\}$ ($1$ for supported, S; $0$ for unsupported, U); depth error
$e=|\hat D_i(p)-D_i(p)|/D_i(p)$. Per-image tokens
$z^{0}_i=E(I_i)$ depend on $I_i$ only; cross-view tokens
$z_i=A(\mathcal{I})$ depend on all views. A score $s$ is any real function;
$\auroc(s)=P(s_S>s_U)+\tfrac12P(s_S=s_U)$, and $\cur_\tau(s)=P(s_U>\tau)$ with
$\tau$ chosen so that $P(s_S>\tau)=0.9$.

% ---------------------------------------------------------------------------
\subsection{Native confidence predicts error, not support}

DUSt3R and its descendants train confidence with a confidence-weighted
regression loss~\cite{wang2024dust3r}
\begin{equation}
  \mathcal{L}(C)=C\,e-\alpha\log C,\qquad \alpha>0,
  \label{eq:confloss}
\end{equation}
applied only where ground-truth depth exists. For backbones whose exact loss
differs, the argument below holds for any loss whose minimiser is a
decreasing function of expected error.

\begin{proposition}[Loss-optimal confidence]\label{prop:conf}
For a head with input $z$, the minimiser of $\E[\mathcal{L}(C)\mid z]$ is
\begin{equation}
  C^\star(z)=\frac{\alpha}{m(z)},\qquad m(z)=\E[e\mid z].
  \label{eq:copt}
\end{equation}
Consequently $C^\star$ carries information about $Y$ only through $m(z)$.
If $e\perp Y\mid z$, then $C^\star$ and $Y$ are conditionally independent
given $m(z)$, and $\auroc(C^\star)$ for S versus U is the \auroc{} of expected
error for separating S from U.
\end{proposition}
\par\noindent\emph{Proof.}~
$\partial_C\E[\mathcal{L}\mid z]=m(z)-\alpha/C=0$ gives Eq.~\eqref{eq:copt};
the second derivative $\alpha/C^2>0$ makes this the minimum. $C^\star$ is a
function of $m(z)$ alone.
\hfill$\square$\par

\noindent\emph{Reading.} The objective contains no term that depends on $Y$, so
nothing pushes the head to encode support even when $z$ contains it. Invented
depth is accurate (\textsc{AbsRel}$_U$ 0.011 to 0.022, main paper, Sec.~4), so
$m(z)$ is nearly the same on S and U, and Eq.~\eqref{eq:copt} predicts high
confidence on invented geometry. This is the leak in main-paper Table~1.
Proposition~\ref{prop:conf} also explains main-paper Gate~11: filtering for
accuracy is an error-ranking task, which $C^\star$ already performs, so adding
a support score does not raise fused-cloud F1. Support and accuracy are
different targets and need different scores.

% ---------------------------------------------------------------------------
\subsection{The matched split forces chance for any single-view score}

\begin{proposition}[Exchangeability]\label{prop:exch}
Suppose that, in the scored view $i$, the pixels of patch A (U) and patch B
(S$_m$) are exchangeable: the joint law of $(I_i,p)$ for $p\!\in\!A$ equals
that for $p\!\in\!B$ (matched area, depth, texture and fill, with roles
assigned independently of view-$i$ appearance). Then every score
$s=f(I_i,p)$, including any function of $z^0_i$, satisfies
\begin{equation}
  \auroc_m(s)=\tfrac12,\qquad \cur_m(s)=0.90 .
  \label{eq:chance}
\end{equation}
\end{proposition}
\par\noindent\emph{Proof.}~
Exchangeability gives $s_A\overset{d}{=}s_B$, so for independent draws
$P(s_B>s_A)=P(s_A>s_B)$ and $\auroc_m=\tfrac12$. With identical laws, the
threshold that passes 90\% of B passes 90\% of A.
\hfill$\square$\par

\begin{proposition}[Information ordering]\label{prop:order}
Let $\mathcal{F}_0=\{f(I_i,p)\}$ and $\mathcal{F}=\{f(\mathcal{I},p)\}$. Since
$\mathcal{F}_0\subset\mathcal{F}$,
$\sup_{\mathcal{F}}\auroc\ge\sup_{\mathcal{F}_0}\auroc$, and the supremum over
each class is attained by the posterior $\eta(x)=P(Y\!=\!1\mid x)$ of its
input~\cite{menon2015corrupted}. Under the matched split, any
$\auroc_m>\tfrac12$ therefore measures information that comes from views
other than $i$.
\end{proposition}

\noindent\emph{Reading.} Eq.~\eqref{eq:chance} gives the reference values for
the shaded columns of main-paper Table~1. Native confidence scores
$\auroc_m$ 0.43 to 0.55 and $\cur_m$ 0.80 to 0.87, close to chance: it acts as
a single-view score even though the backbone's tokens have cross-view inputs,
as Proposition~\ref{prop:conf} predicts. The per-image control is a member of
$\mathcal{F}_0$ by construction, so Eq.~\eqref{eq:chance} explains why it stays
at 0.56 to 0.61. The residual above $\tfrac12$ is the cost of approximate
matching (image-border and placement cues), which the pre-registered 0.60 bound
caps. Proposition~\ref{prop:order} explains why the control is decisive: the
support head and the control share architecture, labels and training, and
differ only in whether their input is in $\mathcal{F}$ or $\mathcal{F}_0$.

% ---------------------------------------------------------------------------
\subsection{Why cross-view tokens carry support}

Support is defined over the view set: $Y(p)=1$ iff the ground-truth point at
$p$ is visible in some view $j\neq i$. It is a property of $\mathcal{I}$, not
of $I_i$. Cross-view attention computes, for each token, a weighting over
tokens in the other views; when the surface at $p$ appears in view $j$, the
token has a matching partner there, and when it does not, it has none. The
backbone must compute this correspondence to triangulate, so the information
is available whether or not any head reads it. A linear probe reaches
\auroc{} 0.99 (main paper, Figure~8). This shows the information is linearly
decodable, which is why a head of 1.7 to 3.3\,M parameters is enough and why a
probe on all four tapped layers is within 0.015 of any head we trained.

\paragraph{Evidence grows with views.}
Model each other view $j$ as contributing an independent log-likelihood-ratio
term $\ell_j$ for $Y$, with mean $\pm\mu/2$ and variance $\sigma^2$. Summing
$k$ views gives separation $d'_k=\sqrt{k}\,\mu/\sigma$, and under a binormal
score model
\begin{equation}
  \auroc_k=\Phi\!\big(\sqrt{k}\,d'_1/\sqrt2\big).
  \label{eq:views}
\end{equation}
From a one-view value of 0.68 to 0.75, Eq.~\eqref{eq:views} predicts 0.75 to 0.83
for two views and 0.79 to 0.88 for three. The measured values are 0.80 to 0.92
for three or more views (main paper, Sec.~5). The per-view term $\mu$ scales with
parallax, which triangulation needs, so a single low-parallax view gives
little evidence. That is consistent with main-paper Gate~9: the one-view fine-tune
stops at 0.71, short of 0.75, because there is little evidence to read, not
because the head is too small.

% ---------------------------------------------------------------------------
\subsection{Learning from the backbone's own reprojection labels}

The label-free head trains on $\tilde Y$, the reprojection test on the
backbone's own depth and cameras. Treat $\tilde Y$ as a corrupted copy of $Y$
with class-conditional flip rates $\rho_{+}=P(\tilde Y\!=\!0\mid Y\!=\!1)$ and
$\rho_{-}=P(\tilde Y\!=\!1\mid Y\!=\!0)$.

\begin{proposition}[Rank preservation~\cite{menon2015corrupted}]\label{prop:noise}
With $\eta(z)=P(Y\!=\!1\mid z)$ and $\tilde\eta(z)=P(\tilde Y\!=\!1\mid z)$,
\begin{equation}
  \tilde\eta(z)=\rho_{-}+(1-\rho_{+}-\rho_{-})\,\eta(z).
  \label{eq:noisy}
\end{equation}
If $\rho_{+}+\rho_{-}<1$, then $\tilde\eta$ is a strictly increasing affine
function of $\eta$ and $\auroc(\tilde\eta)=\auroc(\eta)$. A head that estimates
$\tilde\eta$ by binary cross-entropy has, in the limit, the same \auroc{} as a
head trained on ground truth.
\end{proposition}
\par\noindent\emph{Proof.}~
$\tilde\eta=P(\tilde Y\!=\!1\mid Y\!=\!1)\eta+P(\tilde Y\!=\!1\mid Y\!=\!0)(1-\eta)
=(1-\rho_{+})\eta+\rho_{-}(1-\eta)$. \auroc{} depends only on the ordering of
scores, which a strictly increasing map preserves. Cross-entropy is a proper
loss, so its population minimiser is $\tilde\eta$.
\hfill$\square$\par

\noindent\emph{Why the head can beat its teacher.} The teacher scores each
pixel with its own reprojection, $t=g(\eta(z))+\xi$, where $\xi$ is
per-pixel depth and pose noise that the tokens do not predict. The head
estimates $\E[\tilde Y\mid z]$ and so averages $\xi$ out. Under a binormal
model with signal variance $\sigma^2$ and noise variance $\sigma_\xi^2$, the
teacher's separation drops to $d'/\sqrt{1+\sigma_\xi^2/\sigma^2}$, while the
head keeps $d'$. This matches main-paper Gates~3 and~4: on DTU the label-free
head reaches 0.99 against 0.94 for reprojection.

\noindent\emph{When it fails.} Eq.~\eqref{eq:noisy} needs
$\rho_{+}+\rho_{-}<1$, and the gap shrinks as $(1-\rho_{+}-\rho_{-})$ falls.
In practice $\rho$ also depends on $z$, which can reorder pixels. Both effects
are largest where the teacher's geometry is poor. This is why adapting
VGGT label-free on CO3Dv2 loses 0.09, and why the DTU-trained head is the safe
default (main paper, Sec.~6).

% ---------------------------------------------------------------------------
\subsection{Why labels from fills need the matched control}

Let $F$ be any appearance cue, such as the fill. If training makes $F$ predict
$\tilde Y$, the loss-optimal head uses $F$; this is what happens with the
benchmark's fills (\auroc{} 1.000, with the per-image control learning the
same shortcut). Matched labels make $F$ uninformative by design:
$P(F\mid\tilde Y\!=\!1)=P(F\mid\tilde Y\!=\!0)$, so $I(\tilde Y;F)=0$. The
population minimiser then gains nothing from $F$ and must use what still
predicts $\tilde Y$, which by Proposition~\ref{prop:order} is cross-view
evidence. The per-image control is the test that this worked: it must stay at
chance under the same labels.

% ---------------------------------------------------------------------------
\subsection{The limit: consistency is not observation}

\begin{proposition}[Indistinguishability]\label{prop:limit}
Let world $W_1$ contain a surface $X$ that no view sees, which every view shows
as a patch $P$, and let $W_2$ contain a real surface that looks like $P$ in
every view. If both worlds produce the same images $\mathcal{I}$, then every
score $s=f(\mathcal{I},p)$, and every test built from the backbone's outputs
(which are themselves functions of $\mathcal{I}$), gives the same value in both.
No such score can label the patch pixels U in $W_1$ and S in $W_2$.
\end{proposition}
\par\noindent\emph{Proof.}~
$s$ is a function of $\mathcal{I}$ and $p$, which are equal in the two worlds.
\hfill$\square$\par

\noindent\emph{Reading.} This is main-paper Figure~13 and the first limitation
in Sec.~7. Reprojection, model disagreement and the support head are all
functions of $\mathcal{I}$, so none of them separates a patch that every view
shows. Refilling the hole does separate it, because refilling changes
$\mathcal{I}$: it is an intervention, not a function of the inputs.
Consequently the head certifies multi-view agreement, not truth. Reprojection
consistency is necessary for support, not sufficient, which is why
reprojection is at chance on the patches (0.49, 0.48).

% ---------------------------------------------------------------------------
\subsection{Consistency check between \auroc{} and \cur}

If scores on S and U are Gaussian with equal variance and separation $d'$,
then $\auroc=\Phi(d'/\sqrt2)$, and with $z_{0.9}=\Phi^{-1}(0.9)=1.2816$,
\begin{equation}
  \cur=\Phi\!\big(z_{0.9}-\sqrt2\,\Phi^{-1}(\auroc)\big).
  \label{eq:binormal}
\end{equation}
Eq.~\eqref{eq:binormal} predicts each table's \cur{} from its \auroc{} alone
(Table~\ref{tab:binormal}). For the support head it fits within 0.02 on all
four datasets, and for native confidence on DTU within 0.01. The two metrics
are therefore consistent with each other, not two separate results. The fit is
worst for reprojection on Photo Tourism (0.16), where the score is a count of
views and the Gaussian model does not hold. On the matched split the model
predicts 0.87 to 0.94, above the measured 0.80 to 0.87, because the two
patches' score distributions differ in spread even when their medians agree.

\begin{table}[!t]
\centering\small
\setlength{\tabcolsep}{4pt}
\begin{tabular}{@{}llccc@{}}
\toprule
detector & data & \auroc & \cur{} meas. & \cur{} pred. \\
\midrule
native conf.   & DTU     & 0.82 & 0.49 & 0.50 \\
native conf.   & CO3Dv2  & 0.67 & 0.80 & 0.75 \\
native conf.   & ETH3D   & 0.76 & 0.50 & 0.61 \\
native conf.   & Photot. & 0.82 & 0.46 & 0.50 \\
reprojection   & DTU     & 0.94 & 0.13 & 0.18 \\
reprojection   & ETH3D   & 0.97 & 0.07 & 0.08 \\
reprojection   & Photot. & 0.82 & 0.34 & 0.50 \\
head, label-free & DTU   & 0.99 & 0.03 & 0.02 \\
head, label-free & CO3Dv2& 0.92 & 0.23 & 0.24 \\
head, label-free & ETH3D & 0.97 & 0.09 & 0.08 \\
head, label-free & Photot.& 0.93 & 0.19 & 0.21 \\
\midrule
chance (Prop.~\ref{prop:exch}) & any & 0.50 & --- & 0.90 \\
\bottomrule
\end{tabular}
\caption{Measured \cur{} against the binormal prediction of
Eq.~\eqref{eq:binormal} from \auroc{} alone (values from main-paper
Table~2, averaged over backbones).}
\label{tab:binormal}
\end{table}
% Needs amsmath + amssymb (already loaded by CVPR preambles).
% Add the two entries in refs_additions.bib to refs.bib.
% =====================================================================
\makeatletter
\@ifundefined{c@suppprop}{\newcounter{suppprop}}{}
\renewcommand{\thesuppprop}{S\arabic{suppprop}}
\@ifundefined{proposition}{%
  \newenvironment{proposition}[1][]{\refstepcounter{suppprop}\par\medskip\noindent
    \textbf{Proposition~\thesuppprop}\if\relax\detokenize{#1}\relax\else\ \textbf{(#1)}\fi\textbf{.}\ \itshape}%
   {\par\medskip}}{}
\makeatother
\providecommand{\auroc}{\textsc{auroc}}
\providecommand{\cur}{\textsc{cur}}
\providecommand{\E}{\mathbb{E}}

\section{Why the method works: logical and mathematical justification}
\label{sec:supp-why}

This section states, under explicit assumptions, why native confidence cannot
detect missing support, why cross-view tokens can, why a head trained on its
backbone's own reprojection labels can match a ground-truth head and beat its
teacher, and why no detector of the inputs can catch an invention every view
agrees on. It ends with a check that the reported \auroc{} and \cur{} values
agree with each other.

\paragraph{Notation.}
View $i$ of image set $\mathcal{I}=\{I_1,\dots,I_N\}$; pixel $p$; support label
$Y\!\in\!\{0,1\}$ ($1$ for supported, S; $0$ for unsupported, U); depth error
$e=|\hat D_i(p)-D_i(p)|/D_i(p)$. Per-image tokens
$z^{0}_i=E(I_i)$ depend on $I_i$ only; cross-view tokens
$z_i=A(\mathcal{I})$ depend on all views. A score $s$ is any real function;
$\auroc(s)=P(s_S>s_U)+\tfrac12P(s_S=s_U)$, and $\cur_\tau(s)=P(s_U>\tau)$ with
$\tau$ chosen so that $P(s_S>\tau)=0.9$.

% ---------------------------------------------------------------------------
\subsection{Native confidence predicts error, not support}

DUSt3R and its descendants train confidence with a confidence-weighted
regression loss~\cite{wang2024dust3r}
\begin{equation}
  \mathcal{L}(C)=C\,e-\alpha\log C,\qquad \alpha>0,
  \label{eq:confloss}
\end{equation}
applied only where ground-truth depth exists. For backbones whose exact loss
differs, the argument below holds for any loss whose minimiser is a
decreasing function of expected error.

\begin{proposition}[Loss-optimal confidence]\label{prop:conf}
For a head with input $z$, the minimiser of $\E[\mathcal{L}(C)\mid z]$ is
\begin{equation}
  C^\star(z)=\frac{\alpha}{m(z)},\qquad m(z)=\E[e\mid z].
  \label{eq:copt}
\end{equation}
Consequently $C^\star$ carries information about $Y$ only through $m(z)$.
If $e\perp Y\mid z$, then $C^\star$ and $Y$ are conditionally independent
given $m(z)$, and $\auroc(C^\star)$ for S versus U is the \auroc{} of expected
error for separating S from U.
\end{proposition}
\par\noindent\emph{Proof.}~
$\partial_C\E[\mathcal{L}\mid z]=m(z)-\alpha/C=0$ gives Eq.~\eqref{eq:copt};
the second derivative $\alpha/C^2>0$ makes this the minimum. $C^\star$ is a
function of $m(z)$ alone.
\hfill$\square$\par

\noindent\emph{Reading.} The objective contains no term that depends on $Y$, so
nothing pushes the head to encode support even when $z$ contains it. Invented
depth is accurate (\textsc{AbsRel}$_U$ 0.011 to 0.022, main paper, Sec.~4), so
$m(z)$ is nearly the same on S and U, and Eq.~\eqref{eq:copt} predicts high
confidence on invented geometry. This is the leak in main-paper Table~1.
Proposition~\ref{prop:conf} also explains main-paper Gate~11: filtering for
accuracy is an error-ranking task, which $C^\star$ already performs, so adding
a support score does not raise fused-cloud F1. Support and accuracy are
different targets and need different scores.

% ---------------------------------------------------------------------------
\subsection{The matched split forces chance for any single-view score}

\begin{proposition}[Exchangeability]\label{prop:exch}
Suppose that, in the scored view $i$, the pixels of patch A (U) and patch B
(S$_m$) are exchangeable: the joint law of $(I_i,p)$ for $p\!\in\!A$ equals
that for $p\!\in\!B$ (matched area, depth, texture and fill, with roles
assigned independently of view-$i$ appearance). Then every score
$s=f(I_i,p)$, including any function of $z^0_i$, satisfies
\begin{equation}
  \auroc_m(s)=\tfrac12,\qquad \cur_m(s)=0.90 .
  \label{eq:chance}
\end{equation}
\end{proposition}
\par\noindent\emph{Proof.}~
Exchangeability gives $s_A\overset{d}{=}s_B$, so for independent draws
$P(s_B>s_A)=P(s_A>s_B)$ and $\auroc_m=\tfrac12$. With identical laws, the
threshold that passes 90\% of B passes 90\% of A.
\hfill$\square$\par

\begin{proposition}[Information ordering]\label{prop:order}
Let $\mathcal{F}_0=\{f(I_i,p)\}$ and $\mathcal{F}=\{f(\mathcal{I},p)\}$. Since
$\mathcal{F}_0\subset\mathcal{F}$,
$\sup_{\mathcal{F}}\auroc\ge\sup_{\mathcal{F}_0}\auroc$, and the supremum over
each class is attained by the posterior $\eta(x)=P(Y\!=\!1\mid x)$ of its
input~\cite{menon2015corrupted}. Under the matched split, any
$\auroc_m>\tfrac12$ therefore measures information that comes from views
other than $i$.
\end{proposition}

\noindent\emph{Reading.} Eq.~\eqref{eq:chance} gives the reference values for
the shaded columns of main-paper Table~1. Native confidence scores
$\auroc_m$ 0.43 to 0.55 and $\cur_m$ 0.80 to 0.87, close to chance: it acts as
a single-view score even though the backbone's tokens have cross-view inputs,
as Proposition~\ref{prop:conf} predicts. The per-image control is a member of
$\mathcal{F}_0$ by construction, so Eq.~\eqref{eq:chance} explains why it stays
at 0.56 to 0.61. The residual above $\tfrac12$ is the cost of approximate
matching (image-border and placement cues), which the pre-registered 0.60 bound
caps. Proposition~\ref{prop:order} explains why the control is decisive: the
support head and the control share architecture, labels and training, and
differ only in whether their input is in $\mathcal{F}$ or $\mathcal{F}_0$.

% ---------------------------------------------------------------------------
\subsection{Why cross-view tokens carry support}

Support is defined over the view set: $Y(p)=1$ iff the ground-truth point at
$p$ is visible in some view $j\neq i$. It is a property of $\mathcal{I}$, not
of $I_i$. Cross-view attention computes, for each token, a weighting over
tokens in the other views; when the surface at $p$ appears in view $j$, the
token has a matching partner there, and when it does not, it has none. The
backbone must compute this correspondence to triangulate, so the information
is available whether or not any head reads it. A linear probe reaches
\auroc{} 0.99 (main paper, Figure~8). This shows the information is linearly
decodable, which is why a head of 1.7 to 3.3\,M parameters is enough and why a
probe on all four tapped layers is within 0.015 of any head we trained.

\paragraph{Evidence grows with views.}
Model each other view $j$ as contributing an independent log-likelihood-ratio
term $\ell_j$ for $Y$, with mean $\pm\mu/2$ and variance $\sigma^2$. Summing
$k$ views gives separation $d'_k=\sqrt{k}\,\mu/\sigma$, and under a binormal
score model
\begin{equation}
  \auroc_k=\Phi\!\big(\sqrt{k}\,d'_1/\sqrt2\big).
  \label{eq:views}
\end{equation}
From a one-view value of 0.68 to 0.75, Eq.~\eqref{eq:views} predicts 0.75 to 0.83
for two views and 0.79 to 0.88 for three. The measured values are 0.80 to 0.92
for three or more views (main paper, Sec.~5). The per-view term $\mu$ scales with
parallax, which triangulation needs, so a single low-parallax view gives
little evidence. That is consistent with main-paper Gate~9: the one-view fine-tune
stops at 0.71, short of 0.75, because there is little evidence to read, not
because the head is too small.

% ---------------------------------------------------------------------------
\subsection{Learning from the backbone's own reprojection labels}

The label-free head trains on $\tilde Y$, the reprojection test on the
backbone's own depth and cameras. Treat $\tilde Y$ as a corrupted copy of $Y$
with class-conditional flip rates $\rho_{+}=P(\tilde Y\!=\!0\mid Y\!=\!1)$ and
$\rho_{-}=P(\tilde Y\!=\!1\mid Y\!=\!0)$.

\begin{proposition}[Rank preservation~\cite{menon2015corrupted}]\label{prop:noise}
With $\eta(z)=P(Y\!=\!1\mid z)$ and $\tilde\eta(z)=P(\tilde Y\!=\!1\mid z)$,
\begin{equation}
  \tilde\eta(z)=\rho_{-}+(1-\rho_{+}-\rho_{-})\,\eta(z).
  \label{eq:noisy}
\end{equation}
If $\rho_{+}+\rho_{-}<1$, then $\tilde\eta$ is a strictly increasing affine
function of $\eta$ and $\auroc(\tilde\eta)=\auroc(\eta)$. A head that estimates
$\tilde\eta$ by binary cross-entropy has, in the limit, the same \auroc{} as a
head trained on ground truth.
\end{proposition}
\par\noindent\emph{Proof.}~
$\tilde\eta=P(\tilde Y\!=\!1\mid Y\!=\!1)\eta+P(\tilde Y\!=\!1\mid Y\!=\!0)(1-\eta)
=(1-\rho_{+})\eta+\rho_{-}(1-\eta)$. \auroc{} depends only on the ordering of
scores, which a strictly increasing map preserves. Cross-entropy is a proper
loss, so its population minimiser is $\tilde\eta$.
\hfill$\square$\par

\noindent\emph{Why the head can beat its teacher.} The teacher scores each
pixel with its own reprojection, $t=g(\eta(z))+\xi$, where $\xi$ is
per-pixel depth and pose noise that the tokens do not predict. The head
estimates $\E[\tilde Y\mid z]$ and so averages $\xi$ out. Under a binormal
model with signal variance $\sigma^2$ and noise variance $\sigma_\xi^2$, the
teacher's separation drops to $d'/\sqrt{1+\sigma_\xi^2/\sigma^2}$, while the
head keeps $d'$. This matches main-paper Gates~3 and~4: on DTU the label-free
head reaches 0.99 against 0.94 for reprojection.

\noindent\emph{When it fails.} Eq.~\eqref{eq:noisy} needs
$\rho_{+}+\rho_{-}<1$, and the gap shrinks as $(1-\rho_{+}-\rho_{-})$ falls.
In practice $\rho$ also depends on $z$, which can reorder pixels. Both effects
are largest where the teacher's geometry is poor. This is why adapting
VGGT label-free on CO3Dv2 loses 0.09, and why the DTU-trained head is the safe
default (main paper, Sec.~6).

% ---------------------------------------------------------------------------
\subsection{Why labels from fills need the matched control}

Let $F$ be any appearance cue, such as the fill. If training makes $F$ predict
$\tilde Y$, the loss-optimal head uses $F$; this is what happens with the
benchmark's fills (\auroc{} 1.000, with the per-image control learning the
same shortcut). Matched labels make $F$ uninformative by design:
$P(F\mid\tilde Y\!=\!1)=P(F\mid\tilde Y\!=\!0)$, so $I(\tilde Y;F)=0$. The
population minimiser then gains nothing from $F$ and must use what still
predicts $\tilde Y$, which by Proposition~\ref{prop:order} is cross-view
evidence. The per-image control is the test that this worked: it must stay at
chance under the same labels.

% ---------------------------------------------------------------------------
\subsection{The limit: consistency is not observation}

\begin{proposition}[Indistinguishability]\label{prop:limit}
Let world $W_1$ contain a surface $X$ that no view sees, which every view shows
as a patch $P$, and let $W_2$ contain a real surface that looks like $P$ in
every view. If both worlds produce the same images $\mathcal{I}$, then every
score $s=f(\mathcal{I},p)$, and every test built from the backbone's outputs
(which are themselves functions of $\mathcal{I}$), gives the same value in both.
No such score can label the patch pixels U in $W_1$ and S in $W_2$.
\end{proposition}
\par\noindent\emph{Proof.}~
$s$ is a function of $\mathcal{I}$ and $p$, which are equal in the two worlds.
\hfill$\square$\par

\noindent\emph{Reading.} This is main-paper Figure~13 and the first limitation
in Sec.~7. Reprojection, model disagreement and the support head are all
functions of $\mathcal{I}$, so none of them separates a patch that every view
shows. Refilling the hole does separate it, because refilling changes
$\mathcal{I}$: it is an intervention, not a function of the inputs.
Consequently the head certifies multi-view agreement, not truth. Reprojection
consistency is necessary for support, not sufficient, which is why
reprojection is at chance on the patches (0.49, 0.48).

% ---------------------------------------------------------------------------
\subsection{Consistency check between \auroc{} and \cur}

If scores on S and U are Gaussian with equal variance and separation $d'$,
then $\auroc=\Phi(d'/\sqrt2)$, and with $z_{0.9}=\Phi^{-1}(0.9)=1.2816$,
\begin{equation}
  \cur=\Phi\!\big(z_{0.9}-\sqrt2\,\Phi^{-1}(\auroc)\big).
  \label{eq:binormal}
\end{equation}
Eq.~\eqref{eq:binormal} predicts each table's \cur{} from its \auroc{} alone
(Table~\ref{tab:binormal}). For the support head it fits within 0.02 on all
four datasets, and for native confidence on DTU within 0.01. The two metrics
are therefore consistent with each other, not two separate results. The fit is
worst for reprojection on Photo Tourism (0.16), where the score is a count of
views and the Gaussian model does not hold. On the matched split the model
predicts 0.87 to 0.94, above the measured 0.80 to 0.87, because the two
patches' score distributions differ in spread even when their medians agree.

\begin{table}[!t]
\centering\small
\setlength{\tabcolsep}{4pt}
\begin{tabular}{@{}llccc@{}}
\toprule
detector & data & \auroc & \cur{} meas. & \cur{} pred. \\
\midrule
native conf.   & DTU     & 0.82 & 0.49 & 0.50 \\
native conf.   & CO3Dv2  & 0.67 & 0.80 & 0.75 \\
native conf.   & ETH3D   & 0.76 & 0.50 & 0.61 \\
native conf.   & Photot. & 0.82 & 0.46 & 0.50 \\
reprojection   & DTU     & 0.94 & 0.13 & 0.18 \\
reprojection   & ETH3D   & 0.97 & 0.07 & 0.08 \\
reprojection   & Photot. & 0.82 & 0.34 & 0.50 \\
head, label-free & DTU   & 0.99 & 0.03 & 0.02 \\
head, label-free & CO3Dv2& 0.92 & 0.23 & 0.24 \\
head, label-free & ETH3D & 0.97 & 0.09 & 0.08 \\
head, label-free & Photot.& 0.93 & 0.19 & 0.21 \\
\midrule
chance (Prop.~\ref{prop:exch}) & any & 0.50 & --- & 0.90 \\
\bottomrule
\end{tabular}
\caption{Measured \cur{} against the binormal prediction of
Eq.~\eqref{eq:binormal} from \auroc{} alone (values from main-paper
Table~2, averaged over backbones).}
\label{tab:binormal}
\end{table}